problem:  
Let \((M^n,g)\) be a complete noncompact Riemannian manifold that is **weakly asymptotically conical** (WAC) in the sense that there exists a compact set \(K\subset M\) and a diffeomorphism  
\[
M\setminus K \cong (R_0,\infty)\times S^{n-1}
\]
such that  
\[
g = dr^2 + r^2 h_r, \qquad \|h_r - h_\infty\|_{C^2(S^{n-1})} = O(r^{-\alpha}), \quad \text{for some } 0<\alpha<1.
\]
Unlike standard AC manifolds where \(\alpha>1\), here the decay is too slow for the full Mazzeo–Melrose or Lockhart–McOwen elliptic Fredholm theory to apply.  

Assume \(\mathrm{Ric}_g \ge -K r^{-2}\) for some \(K>0\), and let  
\[
\mathcal{H}^2(M) = \{u\in L^2(M) \mid \Delta_g u=0\}.
\]

**Open problem:**  
For such WAC manifolds:

1. Is the dimension of \(\mathcal{H}^2(M)\) still completely determined by the indicial roots of the Laplacian on the asymptotic cone \(C(S^{n-1},h_\infty)\)? That is, do the same indicial roots governing formal harmonic asymptotics on cones still classify actual \(L^2\)-harmonic modes when the cone perturbation decays only like \(r^{-\alpha}\) with \(\alpha \le 1\)?  

2. Alternatively, can one construct examples of WAC manifolds with the *same* asymptotic cone \((S^{n-1}, h_\infty)\) but *different* \(L^2\)-harmonic dimensions, indicating that the spectral-indicial correspondence breaks down below a certain decay threshold \(\alpha_c\)?  

Equivalently:  
Does there exist a critical rate \(\alpha_c>0\) such that the structure and dimension of the \(L^2\)-harmonic space are invariant under perturbations of order \(O(r^{-\alpha})\) **if and only if** \(\alpha>\alpha_c\)?  

This requires identifying whether indicially controlled Fredholm stability fails for weakly conic ends, and if new \(L^2\)-harmonic states can emerge or vanish under arbitrarily slow geometric convergence to the conical model.

Why is it a "good" problem:  
This formulation isolates the precise frontier where existing elliptic and scattering theories cease to guarantee spectral invariance and vanishing results for harmonic functions. The novelty lies in probing the **threshold of analytic stability** of Liouville-type phenomena under subcritical (\(\alpha\le1\)) decay, beyond the reach of classical asymptotically conic frameworks. It connects core ideas from harmonic analysis, partial differential equations, and geometric asymptotics: how slow geometric convergence impacts the spectral and Fredholm properties of the Laplacian. The problem elegantly bridges geometric analysis and microlocal elliptic theory—simple in statement, but deeply tied to the structure of Sobolev mapping properties, singular continuous spectrum, and vanishing theorems. Resolving it would extend our understanding of when asymptotic geometry “remembers” its spectral-analytic invariants, achieving the “narrow entrance, profound depth” ideal of a good research problem.